\documentclass{beamer}

\usetheme{Antibes}

% Loading preambulo
% PREAMBULO ====================================================================

% Packages ---------------------------------------------------------------------

% Package to Portuguese language
\usepackage[brazilian]{babel}
% Packages to Figures
\usepackage{graphicx}
\usepackage{tikz}
% Packages to math symbols and expressions
\usepackage{amsfonts, amssymb, amsmath}
% Packages to insert code
\usepackage{listings} 
\usepackage{verbatim}
% Package to justify text
\usepackage[document]{ragged2e}
% Package to manage the bibliography
\usepackage[backend=biber, style=numeric, sorting=none]{biblatex}
% Package to facilitate quotations
\usepackage{csquotes}
% Package to use multicols
\usepackage{multicol}
% Packages for url
%	The package libertinus already loads hyperref
%\usepackage[colorlinks=true]{hyperref}
% Font
%   Fira
%\usepackage[sfdefault, lf]{FiraSans}
%   Libertinus
\usepackage{libertinus}
\renewcommand*\familydefault{\sfdefault} %% Only if the base font of the document is to be sans serif
% Colors
\usepackage{xcolor}
\usepackage{colortbl}
% Table
\usepackage{tabularray}
% Fill line
\usepackage{xhfill}

% Configurations ---------------------------------------------------------------

\AtBeginSection{
    \begin{frame}{\secname}
        \tableofcontents[currentsection,hideallsubsections]
    \end{frame}
}

\AtBeginSubsection{
    \begin{frame}{\subsecname}
        \tableofcontents[currentsection,subsectionstyle=show/shaded/hide, subsubsectionstyle=hide]
    \end{frame}
}

\AtBeginSubsubsection{
    \begin{frame}{\subsubsecname}
        \tableofcontents[currentsection,subsectionstyle=show/shaded/hide,subsubsectionstyle=show/shaded/hide/hide]
    \end{frame}
}

% Numbering slides
\setbeamertemplate{footline}[frame number]{}

% Getting rid of bottom navigation bars
\setbeamertemplate{navigation symbols}{}

% Margins
\setbeamersize{text margin left=20pt, text margin right=20pt}

% Colors
\definecolor{burgundy}{RGB}{133,24,42}
\definecolor{bloodred}{RGB}{147,5,5}

\setbeamercolor{structure}{fg=burgundy}

\setbeamercolor*{palette primary}{use=structure,fg=white,bg=structure.fg}
%\setbeamercolor*{palette secondary}{use=structure,fg=white,bg=structure.fg!75!black}
%\setbeamercolor*{palette tertiary}{use=structure,fg=white,bg=structure.fg!50!black}
\setbeamercolor*{palette quaternary}{fg=white,bg=bloodred}

%	hyperref
\hypersetup{
	colorlinks=true,
	urlcolor=blue,
	breaklinks=true}

% New commands -----------------------------------------------------------------

\NewTblrTableCommand \aula{\SetCell{bg=green!60,fg=white}}
\NewTblrTableCommand \prova{\SetCell{bg=red!80,fg=white}}
\NewTblrTableCommand \feriado{\SetCell{bg=blue!50,fg=white}}

\newcommand{\esta}[1]{\textbf{\underline{#1}}}

% Title
\title[GP 05]{Gerência de Projetos}
% Subtitle
\subtitle{Revisão AV1}
% Author of the presentation
\author[E.J.R. Silva]{Evandro J.R. Silva}
% date of the presentation
\date{}

\begin{document}
    
\begin{frame}
    \titlepage
\end{frame}

\begin{frame}{Sumário}
    \tableofcontents
\end{frame}

%===============================================================================
% SECTION 1 ====================================================================
%===============================================================================
\section{Capítulo 1}

\begin{frame}{Capítulo 1}
    \begin{enumerate}
        \justifying
        \item De acordo com o capítulo 1, seção 1.2, do Padrão de Gerenciamento de Projetos, complete: ``O \rule{50pt}{1pt} é um esforço temporário empreendido para criar um produto, serviço ou resultado exclusivo. Já o \rule{50pt}{1pt} é um conjunto de projetos, programas e operações gerenciados como um grupo para atingir objetivos estratégicos.''
        	\begin{enumerate}[a]
        		\item \alert<2>{projeto, portfólio}.
        		\item resultado, programa.
        		\item produto, programa.
        		\item programa, portfólio.
        		\item projeto, resultado.
        	\end{enumerate}
    \end{enumerate}
\end{frame}

\begin{frame}{Capítulo 1}
	\begin{enumerate}
		\setcounter{enumi}{1}
		\justifying
		\item Relacione os termos-chave às suas definições conforme a seção 1.2:
			\begin{columns}
				\begin{column}{0.2\textwidth}
					\begin{enumerate}[A]
						\item Entrega
						\item Resultado
						\item Valor
						\item Produto
					\end{enumerate}
				\end{column}
				
				\begin{column}{0.8\textwidth}
					\begin{itemize}
						\justifying
						\item O benefício que uma organização recebe. É a importância, utilidade ou mérito de algo.
						\item Um item que é produzido, é quantificável e pode ser um item final ou um item intermediário.
						\item Um efeito final ou consequência de um processo ou projeto.
						\item Um artefato, serviço ou resultado tangível ou intangível produzido por um projeto.
					\end{itemize}
				\end{column}
			\end{columns}
			\begin{enumerate}[a]
				\item A, B, C, D.
				\item D, C, B, A.
				\item \alert<2>{C, D, B, A}.
				\item C, A, B, D.
				\item C, D, A, B.
			\end{enumerate}
	\end{enumerate}
\end{frame}

\begin{frame}{Capítulo 1}
	\begin{enumerate}
		\setcounter{enumi}{2}
		\justifying
		\item Analise as afirmativas a seguir e assinale a alternativa correta.
		
		\vspace{10pt}
		``O Guia PMBOK 7ª Edição enfatiza que os projetos existem para entregar resultados (outcomes) e não apenas produtos (outputs).''
		
		\vspace{10pt}
		\centering \textbf{PORQUE}
		
		\vspace{10pt}
		\justifying
		Os resultados permitem que as organizações realizem o valor que o projeto foi planejado para fornecer.
		
		\vspace{10pt}
		\begin{enumerate}[a]
			\item \alert<2>{As duas são verdadeiras e a segunda justifica a primeira.}.
			\item As duas são verdadeiras, mas a segunda não justifica a primeira..
			\item A primeira é verdadeira e a segunda é falsa.
			\item A primeira é falsa e a segunda é verdadeira..
			\item Ambas são falsas.
		\end{enumerate}
	\end{enumerate}
\end{frame}

\begin{frame}{Capítulo 1}
	\begin{enumerate}
		\setcounter{enumi}{3}
		\justifying
		\item Sobre o conceito de Valor (Seção 1.2), analise os itens:
			\begin{enumerate}[I]
				\justifying
				\item O valor é subjetivo e depende de como a organização define o sucesso.
				\item O valor pode ser realizado durante o projeto, no final ou após a conclusão.
				\item O valor é estritamente financeiro, focado em retorno sobre investimento (ROI).
				\item A entrega de valor exige uma mudança de foco de ``entregar requisitos'' para ``entregar benefícios''.
			\end{enumerate}
		\begin{enumerate}[a]
			\item Se apenas I e III estiverem corretas.
			\item Se apenas II e IV estiverem corretas.
			\item Se apenas IV estiver correta.
			\item \alert<2>{Se apenas I, II e IV estiverem corretas}.
			\item Se todas estiverem corretas.
		\end{enumerate}
	\end{enumerate}
\end{frame}

\begin{frame}{Capítulo 1}
	\begin{enumerate}
		\setcounter{enumi}{4}
		\justifying
		\item Segundo o capítulo 2 do Padrão de Gerenciamento de Projetos, podemos inferir que \rule{50pt}{1pt} refere-se ao conjunto de autoridades, políticas, procedimentos e processos que orientam as atividades de gerenciamento de projetos para atingir as metas organizacionais, estratégicas e operacionais.
		\begin{enumerate}[a]
			\item a criação de valor.
			\item as operações.
			\item a organização.
			\item o sistema de entrega de valor.
			\item \alert<2>{a governança}.
		\end{enumerate}
	\end{enumerate}
\end{frame}

\begin{frame}{Capítulo 1}
	\begin{enumerate}
		\setcounter{enumi}{5}
		\justifying
		\item De acordo com a seção 1.2 do Padrão de Gerenciamento de Projetos, o Gerente de Projeto é definido como:
		\begin{enumerate}[a]
			\justifying
			\item O proprietário do produto (\textit{Product Owner}) em todas as abordagens de desenvolvimento.
			\item O responsável final pela aprovação do orçamento e financiamento do projeto.
			\item A autoridade máxima da organização que define a visão estratégica.
			\item \alert<2>{A pessoa designada pela organização para liderar a equipe responsável por atingir os objetivos do projeto}.
			\item O especialista técnico sênior que executa as tarefas mais complexas do cronograma.
		\end{enumerate}
	\end{enumerate}
\end{frame}

\begin{frame}{Capítulo 1}
	\begin{enumerate}
		\setcounter{enumi}{6}
		\justifying
		\item Uma empresa de software está desenvolvendo um novo aplicativo de segurança. Ao final do projeto, eles entregam o código-fonte e o manual (entregas). No entanto, o real objetivo é reduzir em 30\% os ataques cibernéticos aos clientes. No PMBOK 7, essa redução de 30\% é classificada como:
		\begin{enumerate}[a]
			\justifying
			\item Portfólio.
			\item \alert<2>{\textit{Outcome} (Resultado)}.
			\item Ciclo de Vida.
			\item \textit{Output} (Saída).
			\item Governança.
		\end{enumerate}
	\end{enumerate}
\end{frame}

\begin{frame}{Capítulo 1}
	\begin{enumerate}
		\setcounter{enumi}{7}
		\justifying
		\item Analise as seguintes afirmativas sobre a Equipe do Projeto, e então marque a alternativa correta:
		
		\vspace{10pt} \footnotesize
		(\hspace{10pt}) É o conjunto de indivíduos que realizam o trabalho do projeto para atingir seus objetivos.\\
		(\hspace{10pt}) Em ambientes ágeis, a equipe é obrigatoriamente liderada por um gerente de projeto de comando e controle.\\
		(\hspace{10pt}) A equipe pode incluir pessoas de diferentes áreas funcionais ou departamentos.\\
		(\hspace{10pt}) A responsabilidade pela entrega de valor é compartilhada entre os membros da equipe.\\
		(\hspace{10pt}) Consultores externos não podem ser considerados parte da equipe do projeto.
		
		\vspace{10pt} \normalsize
		\begin{enumerate}[a]
			\justifying
			\item V V V V V
			\item \alert<2>{V F V V F}
			\item F F F F F
			\item V F V F V
			\item F V F V F
		\end{enumerate}
	\end{enumerate}
\end{frame}

\begin{frame}{Capítulo 1}
	\begin{enumerate}
		\setcounter{enumi}{8}
		\justifying
		\item Analise as afirmativas a seguir e assinale a alternativa correta.
		
		\vspace{10pt}
		O Padrão de Gerenciamento de Projetos é um documento de referência para os profissionais
		
		\vspace{10pt}
		\centering \textbf{PORQUE}
		
		\vspace{10pt}
		\justifying
		Ele descreve os processos granulares e os passos sequenciais que devem ser seguidos em todos os projetos de infraestrutura.
		
		\vspace{10pt}
		\begin{enumerate}[a]
			\item As duas são verdadeiras e a segunda justifica a primeira.
			\item As duas são verdadeiras, mas a segunda não justifica a primeira.
			\item \alert<2>{A primeira é verdadeira e a segunda é falsa}.
			\item A primeira é falsa e a segunda é verdadeira.
			\item Ambas são falsas.
		\end{enumerate}
	\end{enumerate}
\end{frame}

%===============================================================================
% SECTION 2 ====================================================================
%===============================================================================
\section{Capítulo 2}

\begin{frame}{Capítulo 2}
	\begin{enumerate}
		\justifying
		\item No contexto do PMBOK 7ª Edição, o Sistema para Entrega de Valor é melhor descrito como:
		\begin{enumerate}[a]
			\item Um conjunto de processos lineares que garantem o menor custo possível.
			\item Um software de gestão (PMIS) que automatiza a criação de relatórios de status.
			\item \alert<2>{Uma coleção de atividades de trabalho estratégicas destinadas a construir, sustentar e/ou avançar uma organização}.
			\item O departamento de PMO (Escritório de Projetos) que audita os cronogramas.
			\item A estrutura hierárquica de comando e controle de uma empresa.
		\end{enumerate}
	\end{enumerate}
\end{frame}

\begin{frame}{Capítulo 2}
	\begin{enumerate}
		\setcounter{enumi}{1}
		\justifying
		\item Analise as seguintes afirmativas sobre Criação de valor, e então marque a alternativa correta:
		
		\vspace{10pt} \footnotesize
		(\hspace{10pt}) O valor pode ser gerado através da criação de novos produtos ou serviços.\\
		(\hspace{10pt}) Melhorar a eficiência ou a produtividade é uma forma de criar valor.\\
		(\hspace{10pt}) O valor é estático e, uma vez definido no início do projeto, não deve sofrer alterações.\\
		(\hspace{10pt}) Mudanças legais ou regulatórias podem ser impulsionadores para a criação de valor.\\
		(\hspace{10pt}) A descontinuação de projetos que não estão mais alinhados à estratégia também faz parte do sistema de valor.
		
		\vspace{10pt} \normalsize
		\begin{enumerate}[a]
			\justifying
			\item V V V V V
			\item \alert<2>{V V F V V}
			\item F F F F F
			\item V F V F V
			\item F V F V F
		\end{enumerate}
	\end{enumerate}
\end{frame}

\begin{frame}{Capítulo 2}
	\begin{enumerate}
		\setcounter{enumi}{2}
		\justifying
		\item ``A \rule{50pt}{0.7pt} do projeto fornece a estrutura e as regras para a tomada de decisões e supervisão, garantindo que o projeto permaneça alinhado com as metas da organização e os padrões de conformidade.''
		\begin{enumerate}[a]
			\item criação de valor.
			\item operação.
			\item organização.
			\item escolha.
			\item \alert<2>{governança}.
		\end{enumerate}
	\end{enumerate}
\end{frame}

\begin{frame}{Capítulo 2}
	\begin{enumerate}
		\setcounter{enumi}{3}
		\justifying
		\item De acordo com a seção `Funções Associadas a Projetos', quais das seguintes responsabilidades são comuns à maioria dos projetos, independentemente de quem as exerce?
		
		\begin{enumerate}[I]
			\item Prover supervisão e coordenação.
			\item Apresentar objetivos e feedback.
			\item Facilitar e apoiar o trabalho da equipe.
			\item Realizar o trabalho e contribuir com percepções.
		\end{enumerate}
		
		\vspace{10pt}
		\begin{enumerate}[a]
			\item Apenas I e III estão corretas.
			\item Apenas II e IV estão corretas.
			\item Apenas I, II e IV estão corretas.
			\item Apenas III e IV estão corretas.
			\item \alert<2>{Todas estão corretas}.
		\end{enumerate}
	\end{enumerate}
\end{frame}

\begin{frame}{Capítulo 2}
	\begin{enumerate}
		\setcounter{enumi}{4}
		\justifying
		\item Relacione as funções às suas descrições típicas no Padrão:
		\begin{columns}
			\begin{column}{0.3\textwidth}				
				\begin{enumerate}[A]
					\footnotesize
					\item Patrocinador
					\item Equipe do Projeto
					\item Comitê de Governança
					\item Gerente de Projeto
				\end{enumerate}
			\end{column}
			
			\begin{column}{0.7\textwidth}
				\begin{itemize}
					\justifying
					\footnotesize
					\item Grupo que fornece orientação estratégica e decisões de alto nível.
					\item Pessoa que fornece recursos e apoio, e é responsável pelo sucesso do negócio.
					\item Indivíduos que executam o trabalho técnico para produzir as entregas.
					\item Pessoa responsável por liderar a equipe e coordenar o esforço para atingir os objetivos.
				\end{itemize}
			\end{column}
		\end{columns}
		\begin{enumerate}[a]
			\item A, C, B, D.
			\item A, C, D, B.
			\item \alert<2>{C, A, B, D}.
			\item C, A, D, B.
			\item C, D, B, A.
		\end{enumerate}
	\end{enumerate}
\end{frame}

\begin{frame}{Capítulo 2}
	\begin{enumerate}
		\setcounter{enumi}{5}
		\justifying
		\item Ordene as etapas lógicas de como o valor é geralmente processado no sistema:
		
		(\hspace{10pt}) Realização de benefícios pelas operações ou clientes.\\
		(\hspace{10pt}) Definição da estratégia organizacional.\\
		(\hspace{10pt}) Execução do projeto para criar um produto ou serviço.\\
		(\hspace{10pt}) Identificação de oportunidades e seleção de componentes no portfólio.
		
		\begin{enumerate}[a]
			\item 4, 2, 1, 3.
			\item \alert<2>{4, 1, 3, 2}.
			\item 4, 3, 2, 1.
			\item 3, 4, 1, 2.
			\item 3, 2, 1, 4.
		\end{enumerate}
	\end{enumerate}
\end{frame}

\begin{frame}{Capítulo 2}
	\begin{enumerate}
		\setcounter{enumi}{6}
		\justifying
		\item Um projeto de infraestrutura está enfrentando dificuldades porque as normas ambientais do país mudaram durante a execução. Segundo o PMBOK 7, este é um exemplo de influência do:
		\begin{enumerate}[a]
			\justifying
			\item \alert<2>{Ambiente Externo (Fatores Políticos e Legais)}.
			\item Ambiente Interno (Cultura organizacional).
			\item Sistema de Governança do Projeto.
			\item Portfólio de Tecnologia.
			\item Fluxo de Valor de TI.
		\end{enumerate}
	\end{enumerate}
\end{frame}

\begin{frame}{Capítulo 2}
	\begin{enumerate}
		\setcounter{enumi}{7}
		\justifying
		\item Sobre os Fatores Ambientais Internos, quais itens abaixo se enquadram nesta categoria?
		
		\begin{enumerate}[I]
			\item Cultura, estrutura e governança organizacional.
			\item Distribuição geográfica de instalações e recursos.
			\item Condições de mercado e concorrência.
			\item Recursos de infraestrutura (equipamentos, canais de TI).
		\end{enumerate}
		
		\vspace{10pt}
		\begin{enumerate}[a]
			\item Apenas I e II estão corretas.
			\item Apenas III está correta.
			\item \alert<2>{Apenas I, II e IV estão corretas}.
			\item Apenas I e IV estão corretas.
			\item Todas estão corretas.
		\end{enumerate}
	\end{enumerate}
\end{frame}

\begin{frame}{Capítulo 2}
	\begin{enumerate}
		\setcounter{enumi}{8}
		\justifying
		\item Um gerente de projeto percebe que a equipe está desmotivada devido a conflitos interpessoais constantes. Ele decide atuar como um facilitador, removendo barreiras e incentivando a colaboração. No PMBOK 7, essa função de `facilitar e apoiar' é:
		\begin{enumerate}[a]
			\justifying
			\item Exclusiva do Patrocinador (\textit{Sponsor}).
			\item \alert<2>{Uma função que pode ser exercida pelo gerente de projeto ou por um líder da equipe}.
			\item Uma tarefa administrativa delegada ao RH.
			\item Proibida em abordagens preditivas.
			\item Responsabilidade única do cliente.
		\end{enumerate}
	\end{enumerate}
\end{frame}

\begin{frame}{Capítulo 2}
	\begin{enumerate}
		\setcounter{enumi}{9}
		\justifying
		\item Sobre as funções de 'Prover Supervisão e Coordenação' (Seção 2.2), analise:
		
		\begin{enumerate}[I]
			\item Envolve o planejamento e o monitoramento do trabalho do projeto.
			\item Garante que os resultados do projeto permaneçam alinhados com os objetivos do negócio.
			\item É uma função que deve ser sempre exercida por uma única pessoa física.
			\item Inclui a gestão de riscos e a garantia de qualidade.
		\end{enumerate}
		
		\vspace{10pt}
		\begin{enumerate}[a]
			\item Apenas I e II estão corretas.
			\item \alert<2>{Apenas I, II e IV estão corretas}.
			\item Apenas III e IV estão corretas.
			\item Apenas IV está correta.
			\item Todas estão corretas.
		\end{enumerate}
	\end{enumerate}
\end{frame}

\begin{frame}{Capítulo 2}
	\begin{enumerate}
		\setcounter{enumi}{10}
		\justifying
		\item Analise as seguintes afirmativas sobre Ativos de Processos Organizacionais (APOs), e então marque a alternativa correta:
		
		\vspace{10pt} \footnotesize
		(\hspace{10pt}) Incluem planos, processos, políticas e procedimentos formais e informais.\\
		(\hspace{10pt}) Bases de conhecimento histórico e lições aprendidas são exemplos de APOs.\\
		(\hspace{10pt}) APOs são considerados influências do ambiente externo à organização.\\
		(\hspace{10pt}) A equipe do projeto pode atualizar e adicionar novos APOs ao longo do projeto.\\
		(\hspace{10pt}) Metodologias de gerenciamento de projetos padronizadas pela empresa são APOs.
		
		\vspace{10pt} \normalsize
		\begin{enumerate}[a]
			\justifying
			\item V V V V V
			\item \alert<2>{V V F V V}
			\item F F F F F
			\item V F V F V
			\item F V F V F
		\end{enumerate}
	\end{enumerate}
\end{frame}

\begin{frame}{Capítulo 2}
	\begin{enumerate}
		\setcounter{enumi}{11}
		\justifying
		\item Coloque em ordem a sequência lógica de transformação de dados em valor dentro do sistema de entrega:
		
		(\hspace{10pt}) Geração de Valor (Realização de benefícios).\\
		(\hspace{10pt}) Dados de Desempenho (Coleta de métricas no projeto).\\
		(\hspace{10pt}) Informações para Tomada de Decisão (Análise dos dados).\\
		(\hspace{10pt}) Ajustes Estratégicos (Ações baseadas na análise).
		
		\vspace{10pt}
		\begin{enumerate}[a]
			\item 4, 2, 1, 3.
			\item \alert<2>{4, 1, 2, 3}.
			\item 4, 3, 2, 1.
			\item 3, 4, 1, 2.
			\item 3, 2, 1, 4.
		\end{enumerate}
	\end{enumerate}
\end{frame}

\begin{frame}{Capítulo 2}
	\begin{enumerate}
		\setcounter{enumi}{12}
		\justifying
		\item Qual das opções abaixo representa um Fator Ambiental Externo que pode afetar o projeto?
		\begin{enumerate}[a]
			\justifying
			\item Software de gestão de projetos (ferramentas de agendamento).
			\item Recursos financeiros disponíveis na empresa.
			\item \alert<2>{Considerações éticas e normas culturais da sociedade}.
			\item Capacidades dos funcionários e disponibilidade de pessoal.
			\item Documentação de projetos anteriores (lições aprendidas).
		\end{enumerate}
	\end{enumerate}
\end{frame}

\begin{frame}{Capítulo 2}
	\begin{enumerate}
		\setcounter{enumi}{13}
		\justifying
		\item A função de `Apresentar Objetivos e Feedback' no projeto é tipicamente exercida por aqueles que:
		\begin{enumerate}[a]
			\justifying
			\item Não possuem interesse financeiro no projeto.
			\item \alert<2>{Irão utilizar as entregas ou fornecerão os requisitos para o produto}.
			\item São responsáveis apenas por auditar as finanças.
			\item Executam o trabalho técnico diário.
			\item Estão fora da cadeia de comando da organização.
		\end{enumerate}
	\end{enumerate}
\end{frame}

\begin{frame}{Capítulo 2}
	\begin{enumerate}
		\setcounter{enumi}{14}
		\justifying
		\item Sobre o conceito de Governança Organizacional, analise as afirmações:
		
		\begin{enumerate}[I]
			\item Define a direção estratégica e os parâmetros para o sucesso.
			\item É idêntica à gestão do projeto (\textit{management}).
			\item Garante que os riscos sejam identificados e gerenciados adequadamente.
			\item Fornece o framework para a prestação de contas.
		\end{enumerate}
		
		\vspace{10pt}
		\begin{enumerate}[a]
			\item \alert<2>{Apenas I, III e IV estão corretas}.
			\item Apenas I e II estão corretas.
			\item Apenas II e IV estão corretas.
			\item Apenas IV está correta.
			\item Todas estão corretas.
		\end{enumerate}
	\end{enumerate}
\end{frame}

\begin{frame}{Capítulo 2}
	\begin{enumerate}
		\setcounter{enumi}{15}
		\justifying
		\item Analise as seguintes afirmativas sobre a realização de benefícios (Seção 2.1), e então marque a alternativa correta:
		
		\vspace{10pt} \footnotesize
		(\hspace{10pt}) Benefícios são sempre de natureza tangível e financeira.\\
		(\hspace{10pt}) A realização de benefícios pode ocorrer muito tempo após o encerramento do projeto.\\
		(\hspace{10pt}) O plano de realização de benefícios descreve como e quando o valor será medido.\\
		(\hspace{10pt}) Mudanças no ambiente de negócios podem tornar os benefícios planejados obsoletos.\\
		(\hspace{10pt}) Projetos que não contribuem para a realização de benefícios devem ser acelerados para terminar logo.
		
		\vspace{10pt} \normalsize
		\begin{enumerate}[a]
			\justifying
			\item V V V V V
			\item \alert<2>{F V V V F}
			\item F F F F F
			\item V F V F V
			\item F V F V F
		\end{enumerate}
	\end{enumerate}
\end{frame}

\begin{frame}{Capítulo 2}
	\begin{enumerate}
		\setcounter{enumi}{16}
		\justifying
		\item ``A função de \rule{50pt}{0.7pt} no projeto ajuda a equipe a navegar pela burocracia organizacional, facilita a comunicação com as partes interessadas de alto nível e remove impedimentos sistêmicos.''
		\begin{enumerate}[a]
			\item supervisão e coordenação.
			\item objetivos e \textit{feedbacks}.
			\item \alert<2>{facilitação e apoio}.
			\item realização do trabalho e contribuição com \textit{insights}.
			\item aplicação e especialização.
		\end{enumerate}
	\end{enumerate}
\end{frame}

\begin{frame}{Capítulo 2}
	\begin{enumerate}
		\setcounter{enumi}{17}
		\justifying
		\item Conforme o Padrão, o que acontece com as informações de desempenho geradas por um projeto?
		\begin{enumerate}[a]
			\justifying
			\item Elas são mantidas em sigilo pelo gerente de projeto.
			\item Elas devem ser descartadas assim que o projeto entra em fase de encerramento.
			\item \alert<2>{Elas alimentam o portfólio para apoiar a tomada de decisão sobre investimentos futuros}.
			\item Elas servem apenas para auditar erros da equipe técnica.
			\item Elas são enviadas exclusivamente aos órgãos reguladores governamentais.
		\end{enumerate}
	\end{enumerate}
\end{frame}

\begin{frame}{Capítulo 2}
	\begin{enumerate}
		\setcounter{enumi}{18}
		\justifying
		\item Quais dos itens abaixo são considerados influências do Ambiente Interno relacionadas a TI?
		
		\begin{enumerate}[I]
			\item Sistemas de licenciamento de software.
			\item Servidores e capacidades de armazenamento de dados.
			\item Restrições comerciais de fornecedores externos.
			\item Canais de comunicação estabelecidos na empresa.
		\end{enumerate}
		
		\vspace{10pt}
		\begin{enumerate}[a]
			\item Apenas I e III estão corretas.
			\item Apenas II e IV estão corretas.
			\item \alert<2>{Apenas I, II e IV estão corretas}.
			\item Apenas III está correta.
			\item Todas estão corretas.
		\end{enumerate}
	\end{enumerate}
\end{frame}

\begin{frame}{Capítulo 2}
	\begin{enumerate}
		\setcounter{enumi}{19}
		\justifying
		\item No Guia PMBOK 7, a responsabilidade pela entrega de valor é:
		\begin{enumerate}[a]
			\justifying
			\item Exclusiva do Gerente de Projeto.
			\item Exclusiva do Patrocinador.
			\item \alert<2>{Compartilhada por todos os envolvidos no sistema de entrega de valor}.
			\item Do cliente, no momento em que ele aceita a entrega.
			\item Da equipe técnica, independentemente da estratégia.
		\end{enumerate}
	\end{enumerate}
\end{frame}

\begin{frame}{Capítulo 2}
	\begin{enumerate}
		\setcounter{enumi}{20}
		\justifying
		\item Analise as afirmativas a seguir e assinale a alternativa correta.
		
		\vspace{10pt}
		``O ambiente de um projeto pode mudar rapidamente devido a fatores sociais e políticos externos.''
		
		\vspace{10pt}
		\centering \textbf{PORQUE}
		
		\vspace{10pt}
		\justifying
		Os fatores ambientais externos, como mudanças no clima político ou tendências sociais, estão fora do controle direto da equipe do projeto.
		
		\vspace{10pt}
		\begin{enumerate}[a]
			\item \alert<2>{As duas são verdadeiras e a segunda justifica a primeira}.
			\item As duas são verdadeiras, mas a segunda não justifica a primeira.
			\item A primeira é verdadeira e a segunda é falsa.
			\item A primeira é falsa e a segunda é verdadeira.
			\item Ambas são falsas.
		\end{enumerate}
	\end{enumerate}
\end{frame}

\begin{frame}{Capítulo 2}
	\begin{enumerate}
		\setcounter{enumi}{21}
		\justifying
		\item Analise as seguintes afirmativas sobre as funções da equipe de projeto (Seção 2.2), e então marque a alternativa correta:
		
		\vspace{10pt} \footnotesize
		(\hspace{10pt}) As funções de projeto podem ser preenchidas por uma pessoa ou por um grupo.\\
		(\hspace{10pt}) A função de 'prover liderança' é fixa e imutável durante todo o ciclo de vida.\\
		(\hspace{10pt}) Membros da equipe podem compartilhar responsabilidades de liderança dependendo da tarefa.\\
		(\hspace{10pt}) O Gerente de Projeto é a única função responsável por remover obstáculos.\\
		(\hspace{10pt}) A colaboração da equipe é essencial para a integração de diferentes perspectivas.
		
		\vspace{10pt} \normalsize
		\begin{enumerate}[a]
			\justifying
			\item V V V V V
			\item V V F F V
			\item F F F F F
			\item \alert<2>{V F V F V}
			\item F V F V F
		\end{enumerate}
	\end{enumerate}
\end{frame}

\begin{frame}{Capítulo 2}
	\begin{enumerate}
		\setcounter{enumi}{22}
		\justifying
		\item Analise as afirmativas a seguir e assinale a alternativa correta.
		
		\vspace{10pt}
		``A cultura organizacional é considerada um Fator Ambiental Interno crítico.''
		
		\vspace{10pt}
		\centering \textbf{PORQUE}
		
		\vspace{10pt}
		\justifying
		``la molda as normas de comportamento, a disposição para assumir riscos e a forma como a autoridade é exercida.''
		
		\vspace{10pt}
		\begin{enumerate}[a]
			\item \alert<2>{As duas são verdadeiras e a segunda justifica a primeira}.
			\item As duas são verdadeiras, mas a segunda não justifica a primeira.
			\item A primeira é verdadeira e a segunda é falsa.
			\item A primeira é falsa e a segunda é verdadeira.
			\item Ambas são falsas.
		\end{enumerate}
	\end{enumerate}
\end{frame}

\begin{frame}{Capítulo 2}
	\begin{enumerate}
		\setcounter{enumi}{23}
		\justifying
		\item Em um ambiente de entrega de valor, os \rule{50pt}{1pt} são os resultados finais obtidos, enquanto os \rule{50pt}{1pt} são os produtos ou serviços tangíveis criados.
		\begin{enumerate}[a]
			\item \alert<2>{\textit{outcomes} / \textit{outputs}}.
			\item produtos / portfólios.
			\item projetos / produtos.
			\item produtos / projetos.
			\item sistemas / projetos.
		\end{enumerate}
	\end{enumerate}
\end{frame}

\begin{frame}{Capítulo 2}
	\begin{enumerate}
		\setcounter{enumi}{24}
		\justifying
		\item Coloque em ordem o ciclo de vida da informação de governança (do tático ao estratégico):
		
		(\hspace{10pt}) Relatórios de status da equipe do projeto.\\
		(\hspace{10pt}) Consolidação de desempenho do programa.\\
		(\hspace{10pt}) Tomada de decisão pela governança do portfólio.\\
		(\hspace{10pt}) Ajuste da visão estratégica pela alta administração.
		
		\vspace{10pt}
		\begin{enumerate}[a]
			\item 1, 3, 2, 4.
			\item 1, 2, 4, 3.
			\item 2, 1, 3, 4.
			\item \alert<2>{1, 2, 3, 4}.
			\item 2, 3, 1, 4.
		\end{enumerate}
	\end{enumerate}
\end{frame}

\begin{frame}{Capítulo 2}
	\begin{enumerate}
		\setcounter{enumi}{25}
		\justifying
		\item Um sistema de entrega de valor é considerado dinâmico porque:
		\begin{enumerate}[a]
			\justifying
			\item Ele nunca muda uma vez que o projeto começa.
			\item \alert<2>{Ele deve reagir a mudanças no mercado, tecnologia e ambiente competitivo}.
			\item Ele depende apenas de uma ferramenta de software específica.
			\item Ele funciona de forma isolada do resto da organização.
			\item Ele é focado apenas no encerramento rápido dos projetos.
		\end{enumerate}
	\end{enumerate}
\end{frame}

\begin{frame}{Capítulo 2}
	\begin{enumerate}
		\setcounter{enumi}{26}
		\justifying
		\item Sobre a Governança do Projeto, analise:
		
		\begin{enumerate}[I]
			\item Define quem tem autoridade para tomar decisões sobre mudanças.
			\item Estabelece como o desempenho será medido e reportado.
			\item Garante a conformidade com leis externas e políticas internas.
			\item Substitui a necessidade de liderança por parte do gerente de projeto.
		\end{enumerate}
		
		\vspace{10pt}
		\begin{enumerate}[a]
			\item Apenas I e IV estão corretas.
			\item Apenas II e III estão corretas.
			\item \alert<2>{Apenas I, II e III estão corretas}.
			\item Apenas IV está correta.
			\item Todas estão corretas.
		\end{enumerate}
	\end{enumerate}
\end{frame}

\begin{frame}{Capítulo 2}
	\begin{enumerate}
		\setcounter{enumi}{27}
		\justifying
		\item Analise as seguintes afirmativas sobre as operações dentro do sistema de entrega de valor, e então marque a alternativa correta:
		
		\vspace{10pt} \footnotesize
		(\hspace{10pt}) As operações são esforços temporários com início e fim definidos.\\
		(\hspace{10pt}) Elas sustentam a organização através da execução de atividades contínuas.\\
		(\hspace{10pt}) As entregas dos projetos são frequentemente transferidas para as operações.\\
		(\hspace{10pt}) O feedback operacional não deve influenciar o portfólio de projetos.\\
		(\hspace{10pt}) Operações eficientes ajudam a maximizar o valor dos ativos criados por projetos.
		
		\vspace{10pt} \normalsize
		\begin{enumerate}[a]
			\justifying
			\item V V V V V
			\item V V F F V
			\item F F F F F
			\item \alert<2>{F V V F V}
			\item F V F V F
		\end{enumerate}
	\end{enumerate}
\end{frame}

\begin{frame}{Capítulo 2}
	\begin{enumerate}
		\setcounter{enumi}{28}
		\justifying
		\item Analise as afirmativas a seguir e assinale a alternativa correta.
		
		\vspace{10pt}
		``A gestão de portfólios foca em `fazer os projetos certos'.''
		
		\vspace{10pt}
		\centering \textbf{PORQUE}
		
		\vspace{10pt}
		\justifying
		``O objetivo do portfólio é garantir que os investimentos estejam alinhados com a estratégia organizacional.''
		
		\vspace{10pt}
		\begin{enumerate}[a]
			\item \alert<2>{As duas são verdadeiras e a segunda justifica a primeira}.
			\item As duas são verdadeiras, mas a segunda não justifica a primeira.
			\item A primeira é verdadeira e a segunda é falsa.
			\item A primeira é falsa e a segunda é verdadeira.
			\item Ambas são falsas.
		\end{enumerate}
	\end{enumerate}
\end{frame}

\begin{frame}{Capítulo 2}
	\begin{enumerate}
		\setcounter{enumi}{29}
		\justifying
		\item A função de Realizar o Trabalho e Contribuir com Percepções envolve:
		
		\begin{enumerate}[I]
			\item Fornecer o conhecimento técnico necessário para as entregas.
			\item Colaborar com outros membros da equipe para resolver problemas.
			\item Definir sozinha a data final do projeto sem consultar o gerente.
			\item Sugerir melhorias nos processos de trabalho.
		\end{enumerate}
		
		\vspace{10pt}
		\begin{enumerate}[a]
			\item Apenas I e III estão corretas.
			\item Apenas II e IV estão corretas.
			\item \alert<2>{Apenas I, II e IV estão corretas}.
			\item Apenas IV está correta.
			\item Todas estão corretas.
		\end{enumerate}
	\end{enumerate}
\end{frame}

\begin{frame}{Capítulo 2}
	\begin{enumerate}
		\setcounter{enumi}{30}
		\justifying
		\item Relacione o tipo de influência ao seu exemplo correto:
		\begin{columns}
			\begin{column}{0.3\textwidth}
				\begin{enumerate}[A]
					\small
					\item APO (Base de Conhecimento)
					\item APO (Processos)
					\item Fator Ambiental Interno
					\item Fator Ambiental Externo
				\end{enumerate}
			\end{column}
			
			\begin{column}{0.7\textwidth}
				\begin{itemize}
					\small
					\justifying
					
					
					\item Registro de riscos de projetos similares finalizados.
					\item Procedimentos de segurança e saúde exigidos por lei nacional.
					\item Modelos de planos de gerenciamento de projeto da empresa.
					\item Disponibilidade de recursos financeiros e distribuição geográfica.
				\end{itemize}
			\end{column}
		\end{columns}
		\begin{enumerate}[a]
			\item A, B, C, D.
			\item A, C, D, B.
			\item \alert<2>{A, C, B, D}.
			\item C, A, B, D.
			\item C, D, A, B.
		\end{enumerate}
	\end{enumerate}
\end{frame}

\begin{frame}{Capítulo 2}
	\begin{enumerate}
		\setcounter{enumi}{31}
		\justifying
		\item Em uma estrutura de governança eficaz, a \rule{50pt}{1pt} deve ser clara e os \rule{50pt}{1pt} de desempenho devem ser comunicados de forma transparente.
		\begin{enumerate}[a]
			\item \alert<2>{autoridade (ou prestação de contas) / indicadores}.
			\item responsabilidade / processos.
			\item estrutura organizacional / acordos.
			\item autoridade (ou prestação de contas) / projetos.
			\item estrutura organizacional / indicadores.
		\end{enumerate}
	\end{enumerate}
\end{frame}

\begin{frame}{Capítulo 2}
	\begin{enumerate}
		\setcounter{enumi}{32}
		\justifying
		\item Coloque em ordem o fluxo de valor simplificado:
		
		(\hspace{10pt}) Criação de entregas (Outputs).\\
		(\hspace{10pt}) Necessidade ou Oportunidade de Negócio.\\
		(\hspace{10pt}) Obtenção de resultados e benefícios (Outcomes).\\
		(\hspace{10pt}) Realização de valor para a organização ou sociedade.
		
		\vspace{10pt}
		\begin{enumerate}[a]
			\item 1, 3, 2, 4.
			\item 1, 2, 4, 3.
			\item 2, 1, 4, 3.
			\item \alert<2>{2, 1, 3, 4}.
			\item 2, 3, 1, 4.
		\end{enumerate}
	\end{enumerate}
\end{frame}

\begin{frame}{Capítulo 2}
	\begin{enumerate}
		\setcounter{enumi}{33}
		\justifying
		\item Um Programa no sistema de entrega de valor é definido como:
		\begin{enumerate}[a]
			\justifying
			\item Um projeto único de longa duração.
			\item \alert<2>{Um grupo de projetos e subprogramas relacionados gerenciados de forma coordenada}.
			\item O departamento de suporte administrativo da empresa.
			\item Uma ferramenta de software para controle de tarefas.
			\item O conjunto total de todos os trabalhos da organização.
		\end{enumerate}
	\end{enumerate}
\end{frame}

\begin{frame}{Capítulo 2}
	\begin{enumerate}
		\setcounter{enumi}{34}
		\justifying
		\item Quais itens abaixo são Fatores Ambientais Externos?
		
		\begin{enumerate}[I]
			\item Padrões da indústria e normas técnicas governamentais.
			\item Pesquisas de mercado e tendências de consumo.
			\item Estrutura de liderança e hierarquia da empresa.
			\item Bancos de dados comerciais e informações de benchmarking.
		\end{enumerate}
		
		\vspace{10pt}
		\begin{enumerate}[a]
			\item Apenas I e II estão corretas.
			\item \alert<2>{Apenas I, II e IV estão corretas}.
			\item Apenas III está correta.
			\item Todas estão corretas.
			\item Apenas I e IV estão corretas.
		\end{enumerate}
	\end{enumerate}
\end{frame}

%===============================================================================
% FIM ==========================================================================
%===============================================================================

\begin{frame}
    \centering
    \Large
    FIM
\end{frame}
    
\end{document}
